\documentclass[11pt,letterpaper]{article}

% Essential Packages
\usepackage[margin=1in]{geometry}
\usepackage{amsmath,amssymb}
\usepackage{graphicx}
\usepackage{circuitikz}
\usepackage{tikz}
\usetikzlibrary{shapes,arrows,positioning,calc}
\usepackage{booktabs}
\usepackage{multirow}
\usepackage{xcolor}
\usepackage{fancyhdr}
\usepackage{tcolorbox}
\usepackage{enumitem}
\usepackage{float}
\usepackage{listings}
\usepackage{hyperref}

% Colors
\definecolor{nlcblue}{RGB}{0,51,102}
\definecolor{alertred}{RGB}{204,0,0}
\definecolor{safgreen}{RGB}{0,128,0}

% Header/Footer
\pagestyle{fancy}
\fancyhead[L]{\textbf{NLC STEM Club UAV Project}}
\fancyhead[R]{\textbf{Power Budget Analysis}}
\fancyfoot[C]{\thepage}

% Title
\title{\Huge\textbf{Power Budget Analysis}\\
\Large Electrical System Energy Calculations\\
\large NLC STEM Club - Electrical Systems Team}
\author{February 2026}
\date{}

\begin{document}

\maketitle
\thispagestyle{empty}
\newpage
\section{ELECTRICAL POWER BUDGET ANALYSIS}

\begin{itemize}
  \item Fixed-Wing Environmental Monitoring UAV
  \item NLC STEM Club - Electrical Systems Team
  \item Date: February 20, 2026
\end{itemize}

\section{SYSTEM OVERVIEW}

Primary Power Source: OVONIC 4S LiPo (14.8V nominal, 2200mAh, 130C)\par

Power Conversion:     40A ESC with integrated 5V 3A BEC\par

Available 5V Rail:    5V @ 3A maximum (15W total)\par

\section{SECTION 1: DETAILED CURRENT CONSUMPTION ANALYSIS}

\subsection*{FLIGHT-CRITICAL SYSTEMS (5V Rail)}

Component               Voltage    Idle (mA)   Typical (mA)   Peak (mA)\par

\section{MATEK F405-WMN FC         5V          180          250           350}

BN-220 GPS Module         5V           40           60            80\par

ELRS EP2 Receiver         5V           70           90           120\par

Servos (4x ES08MA II)     5V          ---          ---           ---\par

\begin{itemize}
  \item Single Servo Idle                   8           10            15
  \item Single Servo Moving                ---          180           300
  \item Single Servo Stall                 ---          ---           650
  \item All 4 Servos Idle                  32           40            60
  \item All 4 Servos Moving                ---          720          1200
  \item All 4 Servos (1 stall)             ---          ---          1830
\end{itemize}

SUBTOTAL (Flight Sys)                 322         1120          1750+\par

\subsection*{RESEARCH PAYLOAD SYSTEMS (5V Rail)}

Component               Voltage    Idle (mA)   Typical (mA)   Peak (mA)\par

Arduino Nano (ATmega328)  5V           20           30            50\par

BME280 Sensor (I2C)       3.3V          3            3             4\par

MQ135 Gas Sensor          5V          140          150           165\par

MicroSD Card Module       5V           15           80           200\par

SUBTOTAL (Research)                   178          263           419\par

\section{SECTION 2: POWER CONSUMPTION SCENARIOS}

Scenario                           Current Draw    \% of BEC     Status\par

Ground Idle (All Systems On)         500 mA         16.7\%      $\checkmark$ SAFE\par

Cruise Flight (No Servo Load)        1383 mA        46.1\%      $\checkmark$ SAFE\par

Active Maneuvering (2 servos)        1743 mA        58.1\%      [WARNING] MARGINAL\par

Aggressive Turn (All servos)         2303 mA        76.8\%      [WARNING] HIGH RISK\par

Control Surface Jam (1 stall)        2543 mA        84.8\%      [WARNING] CRITICAL\par

Multiple Servo Stall (2+)            3193+ mA       106\%+      $\times$ OVERLOAD\par

\section{CRITICAL ANALYSIS:}

[WARNING] WARNING: ESC BEC (3A) is INSUFFICIENT for worst-case scenarios\par

The 5V 3A BEC provides 3000mA maximum. Under aggressive maneuvering or\par

servo stall conditions, total system draw exceeds BEC capacity.\par

\section{FAILURE MODES:}

\begin{itemize}
  \item 1. BEC voltage sag $\rightarrow$ Flight controller brownout $\rightarrow$ Loss of control
  \item 2. BEC thermal shutdown $\rightarrow$ Complete power loss $\rightarrow$ Crash
  \item 3. ESC protective cutoff $\rightarrow$ Motor stops mid-flight $\rightarrow$ Crash
\end{itemize}

SECTION 3: MOTOR POWER ANALYSIS (14.8V Rail)\par

Component               Voltage    Current (A)    Power (W)    Notes\par

\section{FLASH HOBBY D2830EVO     14.8V        ---           ---        1000KV}

\begin{itemize}
  \item Idle/No Load                      0.8          11.8W       Propeller friction
  \item 50\% Throttle                      8-12        118-177W     Cruise flight
  \item 75\% Throttle                     15-20        222-296W     Climb/acceleration
  \item Full Throttle                    25-30        370-444W     Maximum power
  \item Max Burst                        35-38        518-562W     Emergency climb
\end{itemize}

ESC Rating: 40A continuous @ 14.8V (4S) = 592W maximum\par

\section{SAFETY MARGIN:}

\begin{itemize}
  \item Cruise:   40A - 12A = 28A remaining (70\% headroom)  $\checkmark$ EXCELLENT
  \item Climb:    40A - 20A = 20A remaining (50\% headroom)  $\checkmark$ GOOD
  \item Full:     40A - 30A = 10A remaining (25\% headroom)  $\checkmark$ ACCEPTABLE
  \item Burst:    40A - 38A = 2A remaining  (5\% headroom)   [WARNING] MARGINAL
\end{itemize}

\section{SECTION 4: BATTERY ANALYSIS}

\subsection{Battery Specifications}

\begin{itemize}
  \item Type:         OVONIC 4S LiPo
  \item Capacity:     2200mAh
  \item Voltage:      14.8V nominal (16.8V fully charged, 14.0V cutoff)
  \item C-Rating:     130C burst, \textasciitilde{}65C continuous
  \item Max Current:  2.2Ah $\times$ 130C = 286A burst (way more than needed)
  \item 2.2Ah $\times$ 65C = 143A continuous
\end{itemize}

\subsection{Battery Capacity Allocation}

\begin{itemize}
  \item Motor power dominates the power budget.
\end{itemize}

\subsection{Flight Time Estimates (Based on Motor Current)}

Throttle    Motor (A)   5V Draw (mA)   Total (mAh)   Flight Time\par

Cruise 50\%    10A         1383mA         10000mA      \textasciitilde{}13.2 minutes\par

Climb 75\%     18A         1743mA         18000mA       \textasciitilde{}7.3 minutes\par

Mixed*        14A         1500mA         14000mA      \textasciitilde{}9.4 minutes\par

\begin{itemize}
  \item Mixed profile: 70\% cruise, 20\% climb, 10\% idle
\end{itemize}

NOTE: 5V system draw is negligible compared to motor (< 1.5A equivalent).\par

\section{SAFE BATTERY CUTOFF:}

\begin{itemize}
  \item Minimum: 3.3V per cell $\times$ 4 = 13.2V
  \item Suggested landing voltage: 3.5V/cell = 14.0V
  \item Reserve capacity: 20\% $\rightarrow$ Land at \textasciitilde{}500mAh remaining
\end{itemize}

\section{SECTION 5: RECOMMENDED POWER ARCHITECTURE}

OPTION A: Single BEC Configuration (NOT RECOMMENDED)\par

\begin{itemize}
  \item ESC BEC (5V 3A) $\rightarrow$ ALL systems
\end{itemize}

\begin{itemize}
  \item Pros:  Simplicity, minimal wiring
  \item Cons:  $\times$ Insufficient current for servo stall conditions
  \item Single point of failure
  \item Risk of brownout $\rightarrow$ flight controller reset $\rightarrow$ crash
  \item No engineering margin
\end{itemize}

\begin{itemize}
  \item VERDICT: UNSAFE for portfolio-level aerospace work
\end{itemize}

OPTION B: Dual-BEC Configuration (RECOMMENDED) $\checkmark$\par

\subsection{Power Rail 1 (ESC BEC):        Power Rail 2 (External BEC)}

\begin{itemize}
  \item Flight Controller            - 4$\times$ Servos
  \item GPS Module                   (High current isolation)
  \item ELRS Receiver
\end{itemize}

\subsection{Power Rail 3 (Arduino Nano onboard regulator from 5V)}

\begin{itemize}
  \item BME280 Sensor
  \item MQ135 Sensor
  \item MicroSD Module
\end{itemize}

\begin{itemize}
  \item Pros:  $\checkmark$ Isolated servo power prevents brownout
  \item Redundancy (if servo BEC fails, FC stays alive)
  \item Engineering margin: \textasciitilde{}1.5A on FC rail
  \item Professional-grade power distribution
  \item Demonstrates systems engineering competency
\end{itemize}

\begin{itemize}
  \item VERDICT: CORRECT ENGINEERING CHOICE for serious aerospace project
\end{itemize}

OPTION C: Hybrid Configuration (ACCEPTABLE ALTERNATIVE)\par

\begin{itemize}
  \item ESC BEC $\rightarrow$ FC + GPS + Receiver + Nano/Sensors
  \item External BEC $\rightarrow$ Servos ONLY
\end{itemize}

\begin{itemize}
  \item Advantage: Maximum servo current headroom
  \item Use if: You want absolute power isolation for control surfaces
\end{itemize}

\section{SECTION 6: POWER DISTRIBUTION IMPLEMENTATION}

\subsection{Required Additional Hardware}

\begin{itemize}
  \item 5V UBEC (5-6A rating recommended)
  \item Example: Hobbywing 5V 6A UBEC or similar
  \item Input: 7-26V (connect to main LiPo via XT60 breakout)
  \item Output: 5V @ 6A (more than enough for 4 servos)
\end{itemize}

\begin{itemize}
  \item XT60 Y-Harness or Distribution Board
  \item To split main battery power to ESC + External BEC
\end{itemize}

\begin{itemize}
  \item Power Distribution PCB (optional but recommended)
  \item Clean solder joints for all 5V connections
  \item Common ground plane
\end{itemize}

\subsection{Wiring Strategy}

\begin{itemize}
  \item 1. Main LiPo $\rightarrow$ XT60 Splitter
  \item +-$\rightarrow$ ESC (motor + 5V BEC output)
  \item +-$\rightarrow$ External BEC (5V output to servo rail)
\end{itemize}

\begin{itemize}
  \item 2. ESC BEC 5V Output $\rightarrow$ FC + GPS + Receiver rail
\end{itemize}

\begin{itemize}
  \item 3. External BEC 5V Output $\rightarrow$ Servo power pins
\end{itemize}

\begin{itemize}
  \item 4. Nano: Powered from FC 5V rail or external BEC
\end{itemize}

\section{5. ALL GROUNDS MUST BE COMMON}

\begin{itemize}
  \item Critical: Connect all ground wires together
\end{itemize}

\section{SECTION 7: RISK ASSESSMENT \& MITIGATION}

Risk \#1: BEC Overload During Flight\par

\begin{itemize}
  \item Likelihood:   HIGH (if using single BEC)
  \item Impact:       CATASTROPHIC (loss of aircraft)
  \item Mitigation:   Use dual-BEC architecture (Option B)
  \item Validation:   Bench test with all servos under load
\end{itemize}

Risk \#2: Electrical Noise from ESC/Motor\par

\begin{itemize}
  \item Likelihood:   MEDIUM
  \item Impact:       MODERATE (sensor errors, GPS dropouts)
  \item Mitigation:   - Twist motor wires
  \item Route motor wires away from signal wires
  \item Use ferrite beads on ESC power lines
  \item Keep MQ135 sensor away from high-current paths
  \item Validation:   Monitor GPS lock during motor operation on bench
\end{itemize}

Risk \#3: Ground Loop / Poor Grounding\par

\begin{itemize}
  \item Likelihood:   MEDIUM
  \item Impact:       MODERATE (erratic behavior, noise)
  \item Mitigation:   - Star ground topology from battery negative
  \item Use thick ground wires (18-20 AWG)
  \item Verify continuity across all grounds
  \item Validation:   Multimeter continuity test
\end{itemize}

Risk \#4: Servo Stall Current Spike\par

\begin{itemize}
  \item Likelihood:   LOW (good airframe should prevent jamming)
  \item Impact:       HIGH (voltage sag, brownout)
  \item Mitigation:   - Dual-BEC architecture
  \item Ensure control surfaces move freely
  \item Add current limiting if needed
  \item Validation:   Manual obstruction test on bench (brief)
\end{itemize}

Risk \#5: LiPo Mishandling\par

\begin{itemize}
  \item Likelihood:   MEDIUM (team inexperience)
  \item Impact:       CATASTROPHIC (fire, injury)
  \item Mitigation:   - Mandatory safety training
  \item LiPo safety bag usage
  \item Low-voltage alarm
  \item Never charge unattended
  \item Validation:   Supervised charging protocol
\end{itemize}

Risk \#6: MQ135 Sensor Warm-Up\par

\begin{itemize}
  \item Likelihood:   CERTAIN (MQ135 requires 24-48hr burn-in)
  \item Impact:       LOW (inaccurate gas readings initially)
  \item Mitigation:   - Power sensor for 24-48 hours before first flight
  \item Document calibration procedure
  \item Log warm-up period in data
  \item Validation:   Monitor sensor readings over 48-hour period
\end{itemize}

\section{SECTION 8: VALIDATION \& TESTING REQUIREMENTS}

PHASE 1: Component-Level Testing\par

\begin{itemize}
  \item Verify all component voltage ratings (multimeter)
  \item Test each servo individually (movement, current draw)
  \item Confirm GPS lock indoors/outdoors
  \item Validate Arduino sensor readings
  \item Measure ESC BEC output under no load: Expected 5.0V $\pm$ 0.1V
\end{itemize}

PHASE 2: Subsystem Integration Testing\par

\begin{itemize}
  \item Power test bench: ESC + BEC + FC + GPS + Receiver
  \item Measure current draw at idle: Expected \textasciitilde{}320mA
  \item Motor spin test (NO PROP): Verify ESC signal response
  \item Servo bench test: All 4 servos moving simultaneously
  \item $\rightarrow$ Measure current with clamp meter: Should NOT exceed 2000mA
  \item Arduino data logging: Confirm MicroSD writes correctly
\end{itemize}

PHASE 3: Full System Integration\par

\begin{itemize}
  \item All systems powered simultaneously
  \item Measure voltage at FC under servo load: Must stay > 4.8V
  \item GPS lock test with motor running: Verify no GPS dropouts
  \item Range test: Verify receiver connection at 100m
  \item Vibration test: Secure all connections
\end{itemize}

PHASE 4: Ground Validation\par

\begin{itemize}
  \item Control surface checks: Full deflection, no binding
  \item Failsafe test: Turn off transmitter $\rightarrow$ verify servo behavior
  \item Battery alarm test: Verify low voltage warning at 14.0V
  \item Emergency shutdown procedure: Kill switch test
\end{itemize}

\section{SECTION 9: CURRENT BUDGET SUMMARY CHART}

\subsection*{5V POWER RAIL BUDGET (ESC BEC: 3000mA Available)}

\subsection*{Flight Controller      ████░░░░░░░░░░░░░░░░  250mA  (8.3%)}

\subsection*{GPS Module             ██░░░░░░░░░░░░░░░░░░   60mA  (2.0%)}

\subsection*{ELRS Receiver          ███░░░░░░░░░░░░░░░░░   90mA  (3.0%)}

\subsection*{Arduino Nano + Sensors ████░░░░░░░░░░░░░░░░  263mA  (8.8%)}

\subsection*{FLIGHT CRITICAL ONLY   █████████████░░░░░░░  663mA  (22.1%) ✓}

\subsection*{Servos (Cruise)        ████████████████████ 720mA  (24.0%)}

\subsection*{TOTAL TYPICAL          ████████████████████████████░ 1383mA (46%)}

\subsection*{Servos (Aggressive)    ████████████████████████████████ 1200mA}

\subsection*{TOTAL PEAK             ████████████████████████████████████ 1863mA}

\subsection*{(62.1% of BEC) ⚠}

\subsection*{MARGIN (Typical)       ███████████████░░░░░░░░░░░░░ 1617mA (54%)}

\subsection*{MARGIN (Peak)          ██████████░░░░░░░░░░░░░░░░░░ 1137mA (38%)}

[WARNING] WARNING: With single BEC, margin drops to 38\% under peak load.\par

\begin{itemize}
  \item Servo stall condition (2543mA) would EXCEED BEC capacity.
\end{itemize}

\begin{itemize}
  \item SOLUTION: Dual-BEC architecture provides 100\% margin on flight-critical
  \item systems, ensuring FC cannot brown out due to servo load.
\end{itemize}

\section{SECTION 10: EXECUTIVE SUMMARY FOR MEETING}

\section{KEY FINDINGS:}

\begin{itemize}
  \item 1. Total 5V system draw: 1383mA typical, 1863mA peak
  \item 2. ESC BEC capacity: 3000mA (3A)
  \item 3. Critical margin concern: Servo stall can exceed BEC limit
  \item 4. Motor power: 10-18A typical (within ESC and battery limits)
  \item 5. Estimated flight time: 9-13 minutes depending on throttle profile
\end{itemize}

\section{RECOMMENDATION:}

\begin{itemize}
  \item Implement DUAL-BEC ARCHITECTURE (Option B) to isolate servo power
  \item from flight-critical systems. This is the only configuration that
  \item provides adequate engineering margin for safe flight operations.
\end{itemize}

\section{IMMEDIATE ACTION ITEMS:}

\begin{itemize}
  \item Purchase additional 5V UBEC (5-6A rating)
  \item Create power distribution board or wiring harness
  \item Implement common ground topology
  \item Begin Phase 1 component testing
\end{itemize}

\section{DOCUMENT REVISION HISTORY}

Version 1.0 | February 20, 2026 | Initial power budget analysis\par

\end{document}
